% Detailed derivation gathered at the end of the owning structure.

\begin{derivation}[从\eqnrefs{eq:group-commutator-bch-dp8}到\eqnrefs{eq:lie-algebra-closure-dp68}]
在量子态 Hilbert 空间中，把单位元附近的连续变换写成

\begin{equation*}
 U(\boldsymbol\theta)
 =\id-\frac{\ii}{\hbar}\theta^aT_a+
 \order(\theta^2).
\end{equation*}
若 $U$ 对实参数保持内积，$U^\dagger U=\id$ 的一阶项给
\begin{equation*}
 T_a^\dagger=T_a,
\end{equation*}
所以这里采用厄米生成元约定。

再取
\begin{equation*}
 U_a(\epsilon)=\ee^{-\ii\epsilon T_a/\hbar},
 \qquad
 U_b(\delta)=\ee^{-\ii\delta T_b/\hbar}.
\end{equation*}
正文\eqnrefs{eq:group-commutator-bch-dp8} 已由 BCH 公式给出
\begin{equation*}
 U_a(\epsilon)U_b(\delta)U_a(-\epsilon)U_b(-\delta)
 =\id-\frac{\epsilon\delta}{\hbar^2}[T_a,T_b]+
 \order(\epsilon^2\delta,\epsilon\delta^2).
\end{equation*}
另一方面，这个乘积仍是群元素，而且在 $\epsilon,\delta\to0$ 时趋近单位元，所以必可写成某个新的无穷小参数 $\delta\theta^c=\order(\epsilon\delta)$：
\begin{equation*}
 U(\delta\boldsymbol\theta)
 =\id-\frac{\ii}{\hbar}\delta\theta^cT_c+
 \order\!\left(\|\delta\boldsymbol\theta\|^2\right).
\end{equation*}
局部光滑性允许把最低阶参数写成
\begin{equation*}
 \delta\theta^c=\epsilon\delta\,f_{ab}{}^c/\hbar,
\end{equation*}
其中 $f_{ab}{}^c$ 与 $\epsilon,\delta$ 无关。比较两个 $\epsilon\delta$ 阶系数便有
\begin{equation*}
 [T_a,T_b]=\ii\hbar f_{ab}{}^cT_c,
\end{equation*}
即正文\eqnrefs{eq:lie-algebra-closure-dp68}。结构常数的反对称性来自对易子的反对称性；Jacobi 恒等式则直接继承自算符对易子的 Jacobi 恒等式。
\end{derivation}
